\chapter{Methodology}
The first paragraph of the methodology chapter is essential for setting the tone and context of your research methods, but it’s best to write it last, after all the sections and subsections of the methodology chapter are completed. By doing so, you'll have a clear understanding of the details of your approach, allowing you to effectively summarize the key points and provide a cohesive introduction. This opening paragraph should briefly outline the research design and methods you've used, without diving into specifics. It should introduce the reader to the general approach (e.g., qualitative, quantitative, or mixed methods) and provide a roadmap for the detailed explanations that follow in the rest of the methodology chapter. Writing it last ensures that the paragraph reflects the full breadth of your methodology and aligns with the structure and content of the entire chapter.

\section{}{Introduction to Methodology}
This section introduces the overall research approach and outlines why the chosen methods are suitable for answering the research question. It sets the context for the subsequent detailed explanations of the methodology.


\section{Research Design}
Here, you describe the overall structure of the research, whether it's qualitative, quantitative, or mixed-methods. You’ll also mention the study's design type (e.g., experimental, descriptive, case study, survey).


\section{Participants or Subjects}
This section details the participants or subjects involved in the study, including how they were selected, sample size, inclusion/exclusion criteria, and any demographic information relevant to the research.


\section{Data Collection Methods}
You describe the specific techniques used to gather data. This could include surveys, interviews, experiments, observations, or archival research. It’s important to justify why these methods were chosen for your research.


\section{Instruments and Tools}
If you used specific tools, equipment, or instruments (e.g., questionnaires, software, or lab apparatus), provide a detailed explanation of how these tools were used to collect data and why they were appropriate.


\section{Data Analysis Methods}
In this section, explain how the data was processed and analyzed. For quantitative studies, this might involve statistical tests, and for qualitative studies, this could include thematic analysis or coding techniques.


\section{Ethical Considerations}
This part covers the ethical aspects of your research, such as how you ensured participants’ confidentiality, informed consent, and any approval from ethics committees or institutions.


\section{Limitations}
Here, you acknowledge any potential limitations or biases in your methodology, such as sample size, data collection challenges, or constraints that might affect the generalizability of the findings.


\section{Timeline and Procedures}
You might also include a section on how the research was conducted over time. This could outline the step-by-step procedures and timeline for conducting the study from start to finish.
